% BANKS-SOURCE: pdf=171 print=161 kind=section id=9.7
\subsection{Renormalization-group equations in dimensional regularization}

% BANKS-SOURCE: pdf=171 print=161 kind=prose
We have seen that one of the peculiarities of dimensional regularization is that
it introduces an extra dimensionful parameter $\mu$ apart from the coefficient of
the quadratic term in the action. We know that this is a good thing, and that it
is appropriate to introduce such a parameter, the \emph{renormalization scale} in
any regularization scheme. The renormalization scale serves as a surrogate for
the cut-off in the renormalized theory. If we recall our discussion of the exact
RGE, general QFTs are defined by unstable flows near a fixed point, and a
renormalization scale was introduced to mark the momentum scale at which the
behavior of a given QFT deviates from its scale-invariant progenitor. We remarked
that it was always possible to exchange the renormalization scale $\mu$ for

% BANKS-SOURCE: pdf=172 print=162 kind=source-page
(a power of) the coefficient of one of the relevant operators in the theory. It is
possible, but not necessary. We can, instead, over-parametrize the theory in
terms of $\mu$ and all of the $G^I$. Then we expect an equation telling us how
the $G^I$ must change with $\mu$ in order to describe the same physical situation.
This is the remnant of the RG flow near the fixed point, and is also called the
RGE of the renormalized theory.

To derive the RGE using dimensional regularization, write
$d_{\mathrm{reg}}=4-\epsilon$ while keeping $D=4$ for the physical theory.
We note that a bare 1PI
Green function (we suppress the momentum dependence of the Green function for
the moment) satisfies
% BANKS-SOURCE: pdf=172 print=162 kind=equation
\[
  \left.\mu\partial_\mu\right|_{\lambda_0}\Gamma_n^0(m_0^2,\lambda_0,\mu)=0
\]
and
% BANKS-SOURCE: pdf=172 print=162 kind=equation
\[
  \Gamma_n^0(m_0^2,\lambda_0,\mu)=Z^{-n/2}\Gamma_n(m^2,\lambda,\mu).
\]
The first of these equations expresses the fact that $\mu$ does not appear in
the bare Green functions, written in terms of $\lambda_0$. We now apply the first
equation to the second, and obtain
% BANKS-SOURCE: pdf=172 print=162 kind=equation
\[
  \left(\mu\partial_\mu+\beta\partial_\lambda
  +\beta_{m^2}\partial_{m^2}-\frac{n}{2}\gamma\right)\Gamma_n=0.
\]
In this equation, the $\mu$ derivative is taken at fixed $\lambda$ and $m^2$,
and we have defined
% BANKS-SOURCE: pdf=172 print=162 kind=equation
\[
  \gamma\equiv\mu\partial_\mu(\ln Z),\qquad
  \beta\equiv(\mu\partial_\mu)_\lambda,qquad
  \beta_{m^2}\equiv(\mu\partial_\mu)_{m^2},
\]
where the derivatives are taken at fixed $\lambda_0,m_0$. These coefficient
functions are finite functions of the renormalized parameters in any
renormalization scheme. Finiteness of $\Gamma_n$ implies that the $\mu$-dependent
piece of $\ln Z,\lambda$, and $m^2$ is finite. Since $\mu$ always appears raised
to the power $k(d_{\mathrm{reg}}-4)$ in $k$th order in the bare perturbation expansion, the
single-pole parts of the expressions for $\lambda,m^2$, and $Z$ in terms of bare
parameters are not $\mu$-dependent. Furthermore, the multiple poles must all
cancel out in the scaling derivatives, which define $\beta_\lambda,\beta_{m^2}$,
and $\gamma$, since these are all finite. In the MS and $\overline{\mathrm{MS}}$
schemes, the renormalization constants have a $\mu$-independent finite part.
Thus, in this scheme, $\beta,\beta_{m^2}$, and $\gamma$ are $\mu$-independent
and come only from the pole part of their definition. By dimensional analysis,
we must have $\beta=\beta(\lambda)$, $\gamma=\gamma(\lambda)$, and
$\beta_{m^2}=m^2\gamma_{m^2}(\lambda)$.

We can solve these equations in the following way. In words, the equation says
that if we rescale $\mu$ we get a term proportional to $\Gamma_n$, as if
$\Gamma_n$ were just rescaling, plus other terms which can be compensated for by
varying the couplings. This corresponds to the idea that the renormalization
scale is playing the role that the cut-off played in the exact RGE. Near the fixed
point, a change in energy scale is compensated for by a change of scale of the
field and a change in the marginal and relevant couplings. So we make an
\emph{Ansatz}:
% BANKS-SOURCE: pdf=172 print=162 kind=equation
\[
  \Gamma_n(\ee^{-s}\mu,\lambda,m^2)
  =\ee^{F(s)}\Gamma_n(\mu,\lambda(s),m^2(s)).
\]

% BANKS-SOURCE: pdf=173 print=163 kind=source-page
On plugging this in, and comparing with the RGE, we find that any function
satisfying this relation will be a solution of the RGE if
% BANKS-SOURCE: pdf=173 print=163 kind=equation
\[
  \frac{\dd F}{\dd s}=\frac{n}{2}\gamma(\lambda(s)),
  \qquad
  \frac{\dd\lambda}{\dd s}=\beta(\lambda(s)),
\]
and
% BANKS-SOURCE: pdf=173 print=163 kind=equation
\[
  \frac{\dd m^2}{\dd s}=m^2\gamma_{m^2}(\lambda(s)).
\]
The boundary conditions on these equations are
% BANKS-SOURCE: pdf=173 print=163 kind=equation
\[
  F(0)=0,\qquad \lambda(0)=\lambda,\qquad m^2(0)=m^2.
\]
So the solution is
% BANKS-SOURCE: pdf=173 print=163 kind=equation
\[
  \Gamma_n(\ee^{-s}\mu,\lambda,m^2)
  =\ee^{\frac{n}{2}\int_0^s\dd u\,\gamma(\lambda(u))}
  \Gamma_n\!\left(\mu,\lambda(s),
  m^2\ee^{\int_0^s\dd u\,\gamma_{m^2}(\lambda(u))\,\dd u}\right).
\]

Now consider the momentum dependence of $\Gamma_n$. In momentum space (and
including the delta function of momentum conservation) $\Gamma_n$ has engineering
dimension $-n$ in mass units. Thus it is of the form $|p|^{-n}$ times a function
that depends on ratios of the dimensionful quantities,
% BANKS-SOURCE: pdf=173 print=163 kind=equation
\[
  \Gamma_n(\mu,\lambda,m^2,\ee^t p_i)
  =\ee^{-nt}\Gamma_n(\ee^{-t}\mu,\lambda,\ee^{-2t}m^2,p_i).
\]
We now use the solution of the RGE to rewrite this as
% BANKS-SOURCE: pdf=173 print=163 kind=equation
\[
  \Gamma_n(\mu,\lambda,m^2,\ee^t p_i)
  =\ee^{-nt}\ee^{\frac{n}{2}\int_0^t\dd u\,\gamma(\lambda(u))}
  \Gamma_n\!\left(\mu,\lambda(t),
  \ee^{-2t}\ee^{\int_0^t\dd u\,\gamma_{m^2}(\lambda(u))\,\dd u}m^2,p_i\right).
\]
We can interpret this equation as saying that, as we change the momentum scale,
the naively dimensionless coupling $\lambda$ ``runs,'' while the renormalized
field and the mass acquire $\lambda(t)$-dependent corrections to their
dimensions. You should understand that engineering dimensions never change,
and that engineering-dimensional analysis is always valid. However, the behavior
of the theory under rescaling of momentum depends on dimensionless ratios of
momenta to the renormalization scale, and can therefore be changed by quantum
corrections.

The fact that $\beta(\lambda)\neq0$ means that our classical assessment of
$\lambda$ as a marginal parameter is wrong. It does flow along RG trajectories,
and its behavior near the free-field fixed point is determined by the leading
term $\beta=b_0\lambda^2+o(\lambda^3)$. In this approximation, the RGE is a
linear equation for $1/\lambda$,
% BANKS-SOURCE: pdf=173 print=163 kind=equation
\[
  \frac{\dd}{\dd t}\frac{1}{\lambda}=-b_0,
\]
whose solution is
% BANKS-SOURCE: pdf=173 print=163 kind=equation
\[
  \lambda(t)=\frac{\lambda}{1-b_0\lambda t}.
\]
We see that, if $b_0$ is positive (recall that $b_0=3/(16\pi^2)$), then, as
momentum gets larger, $\lambda(t)$ increases, whereas as it gets smaller the
coupling decreases. Such a theory is called
